% ---- pg-563 （书页 551） ----

\[
\boldsymbol{\mu}_1=\begin{bmatrix}5.01\\3.43\\1.46\\0.25\end{bmatrix}, \qquad
\boldsymbol{\mu}_2=\begin{bmatrix}6.26\\2.87\\4.91\\1.68\end{bmatrix}
\]

两个估计的协方差矩阵为
\[
\hat{\boldsymbol{\Sigma}}_1=\begin{bmatrix}
0.121\,8 & 0.097\,2 & 0.016\,0 & 0.010\,1\\
0.097\,2 & 0.140\,8 & 0.011\,5 & 0.009\,1\\
0.016\,0 & 0.011\,5 & 0.029\,6 & 0.005\,9\\
0.010\,1 & 0.009\,1 & 0.005\,9 & 0.010\,9
\end{bmatrix}, \qquad
\hat{\boldsymbol{\Sigma}}_2=\begin{bmatrix}
0.453\,0 & 0.120\,9 & 0.448\,9 & 0.165\,5\\
0.120\,9 & 0.109\,6 & 0.141\,4 & 0.079\,2\\
0.448\,9 & 0.141\,4 & 0.674\,8 & 0.285\,8\\
0.165\,5 & 0.079\,2 & 0.285\,8 & 0.178\,6
\end{bmatrix}
\]

本质上, 两种鸢尾被归到一类里. 这一点可以从\hyperref[fig:12.14]{图 12.14} 中看出. \hyperref[fig:12.14]{图 12.14} 显示的是投影到萼片度量的散点图. \hfill$\blacksquare$

\begin{mvfigure}
\mvinc{pg-563}{0.137,0.358,0.865,0.611}
\phantomsection\label{fig:12.14}\mvcap{图 12.14\quad 最佳模型对于萼片度量的散点图}
\end{mvfigure}

\vspace{0.5em}
\section{多维标度变换}
\phantomsection\label{sec:12.6}

本节讨论在低维空间展示(经过变换的)多元数据的方法. 我们在前面曾讨论过类似的问题, 例如对前两个主成分得分作图, 或者对前两个线性判别量的得分作图. 本节中要讨论的方法与上述方法不同. 这里的基本目标是将原始数据“拟合”到一个低维坐标系中, 使得由降维所引起的任何变形达到最小. 变形通常是指原始数据点之间的相似性或相异性(距离). 虽然可以用欧氏距离来度量最终低维结构中点与点的接近程度, 相似性或相异性概念却依赖于它们的定义方法. 我们所提到的那类低维图称为数据的“定位”.

多维标度变换法涉及这样的问题: 当 $N$ 个项目中各对项目之间的相似性(或距离)给定时, 求这些项目在低维空间中的表示, 并使项目间的接近度与原先的相似性(或距离)“大体匹配”.

标度变换也许不可能同原始相似性(距离)的排序严格匹配. 因此, 标度方法试图找到
